\subsection{推理流程}
\label{sec:inference}

训练好的 PMDF-Net 作为前馈去噪模块运作，在推理时无需迭代优化、平均或每试样微调。给定与训练数据相同采样率的长度为$L$的原始或低平均输入信号，网络通过单次前向传递产生相同长度的去噪波形。这与需要$N$次重复采集并逐点求和以实现$10\log_{10}(N)$dB 信噪比(SNR)提升的常规信号平均直接对比。

\textbf{输入和输出格式。}推理输入是单个超声波形——可以是单脉冲采集或$N \in \{1, 2, 4, 8\}$次平均信号，匹配训练输入范围。不需要额外的元数据、参考信号或校准采集。输出是逐点去噪波形$\hat{y} \in \mathbb{R}^L$，可立即用于后续到达时间拾取、幅值提取或缺陷检测流程。

\textbf{计算成本。}推理吞吐量取决于硬件平台和信号长度。
在桌面 GPU（NVIDIA RTX 3080）上，2048 点信号的前向传递约在 2--5 ms 内完成，在批量模式下吞吐量为 200--500 信号/秒。
在中端边缘 GPU（NVIDIA Jetson Orin）上，延迟增至每信号 15--30 ms，足以支持 30--60 信号/秒的实时处理。
还支持仅 CPU 推理路径，用于标准工作站部署，对于相同信号长度可达到 50--100 信号/秒。

\textbf{相对于平均的加速。}考虑实现与 64 次平均等效信噪比(SNR)的实际场景。
PMDF-Net 在 1 次输入上的单次前向传递约在 GPU 上 3--5 ms 内达到 comparable 去噪质量，
而 64 次硬件平均需要 64 倍的采集持续时间加上求和开销。即使计入网络推理，所提方法在有效采集到输出的延迟上提供$10\times$到$100\times$的加速，取决于比较的$N$次平均基线。

\textbf{实时部署。}PMDF-Net 的全卷积架构不包含需要顺序处理的递归连接或时间依赖性，能够在信号间进行 trivial 并行化。
模型量化（INT8）将占用空间降至 50 MB 以下，使 GPU-边缘部署对于原位检测系统切实可行。
将网络应用于新试样时无需微调或域适应：模型从完整铝训练数据泛化到含缺陷铝和 CFRP 目标，无需参数更新，如第~\ref{sec:experiments}节中所验证。